\documentclass[11pt,a4paper]{article}
\usepackage[margin=25mm]{geometry}
\usepackage[T1]{fontenc}
\usepackage{lmodern,microtype}
\usepackage{amsmath,amssymb,amsthm,mathtools,booktabs,array}
\usepackage[hidelinks]{hyperref}
\hypersetup{pdftitle={A spectral-gap classification for flat Schmidt-rank-two chains},pdfauthor={Anonymous},pdfsubject={A complete flat-Schmidt-rank-two classification; unrefereed AI-assisted candidate},pdfkeywords={qutrit, spectral gap, frustration free, finite-size certificate, locally isospectral}}
\usepackage{fancyhdr}
\pagestyle{fancy}\fancyhf{}
\fancyhead[L]{\small Flat Schmidt-rank-two chains}
\fancyhead[R]{\small Evidence Press \textbar\ v1.0.1}
\fancyfoot[C]{\thepage}
\setlength{\headheight}{14pt}
\setlength{\emergencystretch}{2em}
\newtheorem{theorem}{Theorem}[section]
\newtheorem{lemma}[theorem]{Lemma}
\newtheorem{proposition}[theorem]{Proposition}
\newtheorem{corollary}[theorem]{Corollary}
\theoremstyle{remark}\newtheorem{remark}[theorem]{Remark}
\newcommand{\C}{\mathbb C}
\newcommand{\ket}[1]{\lvert #1\rangle}
\newcommand{\bra}[1]{\langle #1\rvert}
\newcommand{\norm}[1]{\left\lVert #1\right\rVert}
\newcommand{\spec}{\operatorname{spec}}
\newcommand{\rank}{\operatorname{rank}}
\newcommand{\ran}{\operatorname{ran}}
\newcommand{\diag}{\operatorname{diag}}
\newcommand{\gap}{\operatorname{gap}}
\newcommand{\Pd}{P_{\mathrm d}}
\newcommand{\Pm}{P_{\mathrm m}}
\title{\textbf{A spectral-gap classification for\\flat Schmidt-rank-two chains}\\[5pt]
\large Qutrit saturation geometry, a complete isospectral fibre,\\and certified extensions}
\author{Anonymous}
\date{28 September 2026 \quad Version 1.0.1}
\begin{document}
\maketitle
\begin{center}\small Version DOI: \url{https://doi.org/10.5281/zenodo.23024569}\end{center}
\begin{abstract}
We classify the open-chain spectral gap for every rank-one nearest-neighbour interaction whose forbidden vector has exactly two equal nonzero Schmidt probabilities, in any finite local dimension. The chain is gapless precisely when the vector is a balanced marker $(u\otimes w-w\otimes v)/\sqrt2$, with $w$ orthogonal to $u$ and $v$. Every gapless case has the exact gap $1-\cos(\pi/N)$. Otherwise the chain is uniformly gapped. When the two Schmidt supports intersect in a line, a four-site saturation argument gives the explicit bound $3\tau(2-t)(1-t)/1024$, where $t$ is their nontrivial principal-angle cosine and $\tau=2|\langle w,w|\psi\rangle|^2$. We also classify the complete qutrit fibre with flat rank-two Schmidt data and three-site spectrum $0^{(21)},(1/2),1^{(4)},(3/2)$. It has a one-parameter on-site-unitary normal form $(\cos\theta\,00+\sin\theta\,01-12)/\sqrt2$ and is gapped for every $\theta<\pi/2$, with gap at least $\cos^2\theta/6$. The endpoint is gapless although its two- and three-site spectra and every finite-length ground-space dimension agree with those of the gapped members. We retain an exact full-complex rank-one neighbourhood certificate and an opposing-bias rank-two exponential gap upper bound as complementary results. All infinite-length conclusions follow from analytic proofs. Exact finite certificates and independent-implementation numerical diagnostics accompany the manuscript. The classification does not extend to arbitrary Schmidt data or arbitrary rank-two qutrit projectors.
\end{abstract}
\noindent\textbf{Status.} AI-assisted, unrefereed research candidate. Version 1.0.1 addresses a supplied version-1.0 review: the executable classifier now rejects unspecialised symbolic parameters; antecedent locators and the full Gram decomposition are explicit. The review's external independence has not been authenticated. No proof-assistant formalisation, external endorsement, historical-priority guarantee or official REF rating is claimed.

\section{Model, scope and contribution}\label{sec:scope}
For a normalised vector $\psi\in\C^d\otimes\C^d$, let $P=\ket\psi\bra\psi$ and
\begin{equation}\label{eq:HN}
 H_N(P)=\sum_{i=1}^{N-1}P_{i,i+1},\qquad
 \gamma_N(P)=\min\bigl(\spec H_N(P)\setminus\{0\}\bigr),\quad N\ge2.
\end{equation}
The boundary is open and there are no additional boundary terms. The same ordered interaction acts on every bond. The normalisation is $P^2=P$. A uniform gap means $\inf_{N\ge2}\gamma_N>0$; the gapless cases established here satisfy $\gamma_N\to0$. This is a classification by the open-chain spectral gap, not a classification of phases by symmetry or spectral-flow equivalence. Local projector rank and the Schmidt rank of its forbidden vector are different quantities.

The main result classifies \emph{all} rank-one interactions whose forbidden vector has two equal nonzero Schmidt probabilities. It applies in every finite local dimension. The new case is the three-dimensional span of two distinct, intersecting Schmidt supports. A four-site argument distinguishes a propagating extremal overlap mode from an obstruction to its propagation. The result also classifies the complete qutrit fibre with Schmidt probabilities $(1/2,1/2,0)$ and the specified two- and three-site spectra. That fibre has a one-parameter normal form, is uniformly gapped except at its marker endpoint, and has a quantitative lower bound throughout its gapped part.

Bravyi and Gosset classified arbitrary frustration-free qubit chains and proposed rank-one and rank-two qutrit interactions as a next case \cite{BG}. Our theorem does not classify arbitrary qutrit interactions: unequal Schmidt probabilities and general rank-two projectors remain outside its principal scope. Lemm already supplied deterministic three-site criteria and deterministic neighbourhoods \cite{Lemm}; Hunter-Jones and Lemm extend low-rank methods to bounded-degree graphs \cite{HJL}. The hierarchy of Rai et al. includes established finite-size positive decompositions \cite{Rai}. We use that general method, with an explicitly proved open-boundary decomposition, and do not claim a new semidefinite hierarchy.

The previous version supplied examples and bounds but left the isospectral interpolation unresolved. The supplied review correctly identified a weaker locally-isospectral obstruction obtainable from the qubit classification and missing Motzkin comparisons. Section~\ref{sec:qubitcomparison} includes that comparison; Section~\ref{sec:motzkin} gives the actual local terms, normalisations and boundaries. The unrestricted obstruction is not our originality claim. The principal candidate contribution is now the complete flat-Schmidt-rank-two classification and its quantitative four-site proof. A targeted primary-source search did not locate this theorem; that search is not proof of historical priority.

\begin{lemma}[Frustration freeness]\label{lem:ff}
If a two-site projector has rank strictly below $d$, every open chain has a nonzero product vector in its kernel.
\end{lemma}
\begin{proof}
Choose a nonzero first-site vector. Once $x_i$ has been chosen, orthogonality of $x_i\otimes x_{i+1}$ to the forbidden space imposes fewer than $d$ homogeneous linear constraints on $x_{i+1}$. A nonzero solution exists. Continue along the chain. This is the elementary product-state argument underlying the low-rank regime discussed in \cite{Mov}.
\end{proof}

\paragraph{Known components and candidate extensions.}
The qubit dichotomy is Theorem~1 of Bravyi--Gosset~\cite{BG}.
The ground-space recurrence is equation~(11), Section~II, of
Movassagh et al.~\cite{Mov}; our Proposition~\ref{prop:degeneracy}
proves its exact persistence throughout a particular nongeneric fibre,
including the gapless endpoint. The unbiased hopping gap is used in Step~1
of the Motzkin proof~\cite{Motzkin}; Theorem~\ref{thm:marker} permits
arbitrary overlap of the two background vectors and proves a global,
rather than sector-only, equality. The principal classification and the
quantitative saturation obstruction are distinguished from these ingredients.

\section{Complete classification at flat Schmidt rank two}\label{sec:main}
Let
\[
 \rho_L=\operatorname{Tr}_2 P,\qquad \rho_R=\operatorname{Tr}_1 P,
 \qquad L=\ran\rho_L,\quad R=\ran\rho_R.
\]
Both supports are subspaces of the same on-site space, in the same physical basis. Throughout this section the two nonzero eigenvalues of each reduced density matrix are $1/2$. Thus $\Pi_L=2\rho_L$ and $\Pi_R=2\rho_R$ are rank-two orthogonal projectors.

\begin{theorem}[Flat-Schmidt-rank-two dichotomy]\label{thm:classification}
For every finite $d\ge2$, the chain \eqref{eq:HN} is gapless if and only if there are unit vectors $u,v,w\in\C^d$, with $w\perp u,v$, for which, up to an irrelevant overall phase,
\begin{equation}\label{eq:balancedmarker}
 \psi=\frac{\ket{u,w}-\ket{w,v}}{\sqrt2}.
\end{equation}
In every gapless case the full finite-chain gap is exactly
\begin{equation}\label{eq:diffusive}
 \gamma_N=1-\cos(\pi/N).
\end{equation}
Every other interaction in the specified Schmidt class is uniformly gapped. More explicitly:
\begin{enumerate}
\item If $L\cap R=\{0\}$, then $\gamma_N\ge1-\norm{\Pi_L\Pi_R}>0$.
\item If $L=R$, then \eqref{eq:balancedmarker} holds and \eqref{eq:diffusive} applies.
\item If $L\cap R=\C w$, choose unit $u\in L\ominus\C w$ and $v\in R\ominus\C w$, and define
\begin{equation}\label{eq:invariants}
 t=|\langle u,v\rangle|=\sqrt{\operatorname{tr}(\Pi_L\Pi_R)-1}<1,
 \qquad \tau=2|\langle w,w|\psi\rangle|^2\in[0,1].
\end{equation}
The chain is gapless exactly when $\tau=0$. For $\tau>0$, every $N\ge2$ satisfies
\begin{equation}\label{eq:generalbound}
 \gamma_N\ge \frac{3\tau(2-t)(1-t)}{1024}>0.
\end{equation}
\end{enumerate}
\end{theorem}
The invariants do not depend on phases chosen for $u,v,w$. In local dimension three, only the second and third alternatives occur. The bound in \eqref{eq:generalbound} is deliberately conservative; it is not asserted to be the thermodynamic gap or to give an optimal critical exponent.

\begin{corollary}[Dimension-minimal flat separation]\label{cor:minimal}
Every maximally entangled two-qubit forbidden vector gives the exact gap \eqref{eq:diffusive}. Dimension three is minimal for a gapped/gapless separation with flat rank-two Schmidt probabilities, even when complete two- and three-site spectra are also required to agree.
\end{corollary}
The existence part is furnished by Section~\ref{sec:isospectral}. The qubit impossibility follows from the coincident-support alternative of Theorem~\ref{thm:classification}; it also follows immediately from Bravyi--Gosset's transfer-matrix criterion.

\begin{remark}[Higher flat Schmidt ranks]
For a rank-one projector whose forbidden vector has $r\ge3$ equal nonzero Schmidt probabilities, the elementary overlap bound gives $\gamma_N\ge1-2/r>0$. The threshold case $r=2$ is therefore the only flat entangled Schmidt class requiring the saturation analysis below. This observation is a consequence of the standard overlap method, not a separate novelty claim.
\end{remark}

\section{Overlap bounds and an exactly solvable marker family}
The following is the standard adjacent-projection strategy, retaining a finite-path cosine factor. We include the proof and make no claim that this general gap method is new; compare \cite{Lemm,HJL}.
\begin{lemma}[Finite-path overlap bound]\label{lem:overlap}
Let $P$ be a two-site projector and let $c=\norm{P_{12}P_{23}}$. Whenever the right-hand side is positive,
\begin{equation}\label{eq:pathbound}
 \gamma_N(P)\ge1-2c\cos(\pi/N).
\end{equation}
For $P=\ket\psi\bra\psi$, write $\psi=\sum_{i,j}M_{ij}\ket{ij}$ with $\norm M_F=1$. Then
\begin{equation}\label{eq:contraction}
 c=\norm{\overline M M}=\norm{M\overline M}.
\end{equation}
\end{lemma}
\begin{proof}
Put $p_i=P_{i,i+1}$ and $z_i=\norm{p_i x}$ for an arbitrary vector $x$. Disjoint projections commute, hence have nonnegative anticommutator. For adjacent terms, the norm of their product bounds the angle between their ranges, giving
\[
 \langle x,H_N^2x\rangle\ge\sum_{i=1}^{N-1}z_i^2
       -2c\sum_{i=1}^{N-2}z_i z_{i+1}.
\]
The largest eigenvalue of the adjacency matrix of the $(N-1)$-vertex path is $2\cos(\pi/N)$. Thus $H_N^2\ge(1-2c\cos(\pi/N))H_N$, which implies \eqref{eq:pathbound} by the spectral theorem.

For \eqref{eq:contraction}, use the isometries $Vx=\psi\otimes x$ and $Wy=y\otimes\psi$ into three sites. Their overlap has entries
$(V^*W)_{c,a}=\sum_b\overline{M_{ab}}M_{bc}$, so
$V^*W=(\overline M M)^T$. As $P_{12}P_{23}=V(V^*W)W^*$, its norm is the norm of this $d\times d$ overlap matrix. Complex conjugation gives the second equality.
\end{proof}
The readily computable region $\norm{\overline M M}<1/2$ therefore has a uniform gap. It is sufficient, not necessary. Even the largest Schmidt probability being below $1/2$ is sufficient, since $\norm{\overline M M}\le\norm M^2$, but the flat rank-two case reaches the threshold and requires more information.

\begin{theorem}[Exact marker family]\label{thm:marker}
Let $u,v,w$ be unit vectors with $w\perp u,v$. The overlap $\langle u,v\rangle$ is arbitrary. For $a,b\ge0$ with $a^2+b^2=1$, set
\begin{equation}\label{eq:marker-family}
 \psi=a\ket{u,w}-b\ket{w,v}.
\end{equation}
Then for all $N\ge2$,
\begin{equation}\label{eq:marker-exact}
 \gamma_N(\ket\psi\bra\psi)=1-2ab\cos(\pi/N).
\end{equation}
\end{theorem}
\begin{proof}
Choose coordinates with $w$ a real basis vector. The coefficient matrix is
$M=a u w^T-b w v^T$, so
\[
 M\overline M=-ab\bigl(u v^*+\langle u,v\rangle w w^*\bigr).
\]
The two displayed summands have orthogonal domain and range supports; the first has norm one and the second norm at most one. Hence $c=ab$, and Lemma~\ref{lem:overlap} gives the lower bound.

For the matching upper bound, consider the $N$ orthonormal vectors
\[
 e_j=u^{\otimes(j-1)}\otimes w\otimes v^{\otimes(N-j)},\quad 1\le j\le N.
\]
They are orthogonal because the unique $w$ occurs at different sites, regardless of $\langle u,v\rangle$. Their span is reducing: away from the marker, every bond is orthogonal to the forbidden vector, while a bond adjoining it moves the marker by one site. The restricted Hamiltonian is $D^*D$, where row $j$ of $D$ has entries $-b$ at column $j$ and $a$ at column $j+1$. The matrix $DD^*$ has diagonal one and off-diagonal entries $-ab$. Its eigenvalues are
$1-2ab\cos(k\pi/N)$, $1\le k\le N-1$. The least of these attains the lower bound. The cases $ab=0$ also follow directly.
\end{proof}
When $u,v$ are independent and $ab>0$, their union with $w$ spans three local dimensions. Thus the statement includes genuinely qutrit forbidden states, not only a qubit interaction with an unused level. Taking $u=\ket0$, $w=\ket1$, $v=\ket2$, $a=b=1/\sqrt2$ gives the gapless endpoint in Section~\ref{sec:isospectral}.

\section{Four-site saturation and proof of the dichotomy}\label{sec:saturation}
\subsection{Schmidt-support geometry}
Write $\psi=\sum M_{ij}\ket{ij}$. A Schmidt decomposition gives $M=AB^T/\sqrt2$, where $A,B$ are $d\times2$ isometries onto $L,R$. The nonzero singular values of $(\overline M M)^T$ are one half the singular values of $A^*B$. Indeed, the outer factors in $\overline A(B^*A)B^T/2$ are isometries or coisometries. Therefore
\begin{equation}\label{eq:angleoverlap}
 \norm{P_{12}P_{23}}=\tfrac12\norm{\Pi_L\Pi_R}.
\end{equation}
If $L\cap R=\{0\}$, compactness of unit spheres gives $\norm{\Pi_L\Pi_R}<1$ and Lemma~\ref{lem:overlap} proves alternative 1.

Suppose henceforth that $L\cap R=\C w$. Choose phases so $\langle u,v\rangle=t\ge0$. In the orthonormal left frame $(u,w)$ and right frame $(w,v)$,
\begin{equation}\label{eq:Uframe}
 \psi=\frac{a\ket{u,w}+b\ket{u,v}+c\ket{w,w}+d\ket{w,v}}{\sqrt2},
 \qquad U=\begin{pmatrix}a&b\\c&d\end{pmatrix}\in U(2).
\end{equation}
Here $\tau=|c|^2$. These coefficients are unrelated to the weights of the marker or dimer families. Put
\begin{equation}\label{eq:lr}
 \ell=\sqrt2(I\otimes\bra w)\psi=a u+c w,
 \qquad r=\sqrt2(\bra w\otimes I)\psi=c w+d v.
\end{equation}
Both vectors are unit. Unitarity of $U$ gives the useful scalar identities
\begin{equation}\label{eq:scalars}
 \langle\psi|\ell,\ell\rangle=\langle\psi|r,r\rangle
 =\frac c{\sqrt2}=\langle w,w|\psi\rangle.
\end{equation}
For example, the first contraction is
$[c(|a|^2+|c|^2)+ta(\overline b a+\overline d c)]/\sqrt2=c/\sqrt2$.
The second follows from $\overline a c+\overline b d=0$ and $|c|^2+|d|^2=1$. This calculation permits arbitrary complex coefficients; it is not restricted to real forbidden vectors.

\subsection{The extremal adjacent-pair mode}
For adjacent bond projectors $p,q$, define
\begin{equation}\label{eq:Adef}
 \mathcal A(p,q)=(p+q)^2-\tfrac12(p+q).
\end{equation}
By \eqref{eq:angleoverlap}, $\norm{pq}=1/2$. The nonzero singular values of the range overlap are $1/2,t/2$. The two-projection decomposition thus gives eigenvalues $1\pm1/2$ and $1\pm t/2$, with any remaining positive eigenvalues equal to one, for $p+q$. Consequently $\mathcal A(p,q)\ge0$.

With $Vx=\psi\otimes x$ and $Wy=y\otimes\psi$, one has $V^*W\ell=r/2$ and $W^*Vr=\ell/2$. Thus
\begin{equation}\label{eq:xi}
 \xi=\psi\otimes r-\ell\otimes\psi,
 \quad \norm\xi=1,
 \quad p\xi=\tfrac12\psi\otimes r,
 \quad q\xi=-\tfrac12\ell\otimes\psi.
\end{equation}
Since $t<1$, the eigenvalue $1/2$ of $p+q$ is simple. The kernel of $\mathcal A(p,q)$ is exactly the direct sum of the common zero space of $p,q$ and $\C\xi$. On its orthogonal complement,
\begin{equation}\label{eq:delta}
 \mathcal A(p,q)\ge \delta I,
 \qquad \delta=(1-t/2)(1-t)/2=\frac{(2-t)(1-t)}4.
\end{equation}
The statement remains valid after tensoring with arbitrary exterior sites.

\subsection{A clipped four-site criterion}
\begin{lemma}[Four-site saturation criterion]\label{lem:saturation}
Let $P$ be a two-site projector with $\norm{P_{12}P_{23}}\le1/2$. On four sites put $p_j=P_{j,j+1}$, $h=p_1+p_2+p_3$, and
\begin{equation}\label{eq:Gdef}
 \mathcal G=\mathcal A(p_1,p_2)+\mathcal A(p_2,p_3)+2p_1p_3
 =h^2-\tfrac12(p_1+p_3)\ge0.
\end{equation}
If $\mathcal G\ge\kappa h$ for some $0<\kappa\le(3-\sqrt5)/4$, then every open chain satisfies
\begin{equation}\label{eq:saturationbound}
 H_N^2\ge\tfrac32\kappa H_N,\qquad \gamma_N\ge\tfrac32\kappa.
\end{equation}
\end{lemma}
\begin{proof}
Extend $p_i=0$ outside $1\le i\le N-1$ and sum \eqref{eq:Gdef} over all three-bond windows meeting the chain. A clipped one-bond window gives at least $p/2$. A clipped two-bond window gives $(p+q)^2-p/2$, or its reflected version. For any $x$, writing $u=\norm{px}$ and $v=\norm{qx}$ bounds this quadratic form below by $u^2/2+v^2-uv$. Its least matrix eigenvalue is $\eta=(3-\sqrt5)/4$, so it is at least $\eta\langle x,(p+q)x\rangle$. Thus every full or clipped window obeys the same inequality with $\kappa$.

Every bond is counted three times on the right. Direct expansion on the left gives
\[
 \sum_j\mathcal G_j
 =2H_N+2\sum_i\{p_i,p_{i+1}\}+2\sum_i p_i p_{i+2}
 \le2H_N^2.
\]
The difference consists of positive products of disjoint commuting projectors. Therefore $3\kappa H_N\le2H_N^2$. The argument includes $N=2,3$ because the zero extension retains all clipped windows.
\end{proof}
This is an explicit local positive decomposition in the established finite-size framework. Its role is to exploit \emph{incompatible saturation} of adjacent pairs, not to lower the generic three-site threshold by assertion.

\subsection{Quantitative incompatibility of saturation}
\begin{lemma}[The overlap obstruction]\label{lem:quantitative}
Under \eqref{eq:Uframe}, if $t<1$ and $\tau>0$, then
\begin{equation}\label{eq:Gquant}
 \mathcal G\ge\frac{\tau\delta}{128}\,h.
\end{equation}
\end{lemma}
\begin{proof}
For arbitrary $x$ on four sites let $E=\langle x,\mathcal Gx\rangle$ and
\[
 \epsilon^2=\frac{\langle x,\mathcal A(p_1,p_2)x\rangle}{\delta},\quad
 \zeta^2=\frac{\langle x,\mathcal A(p_2,p_3)x\rangle}{\delta},\quad
 R_0=\sqrt{E/\delta}.
\]
Then $\epsilon^2+\zeta^2\le R_0^2$. Orthogonally project $x$ onto the kernels of these two pair operators and use \eqref{eq:xi}. The components in the common zero spaces disappear when the corresponding bond projectors act. Absorbing the factor $1/2$ in exterior vectors $z,y$, we obtain
\begin{align*}
 p_1x&=\psi\otimes r\otimes z+e_1,&
 p_2x&=-\ell\otimes\psi\otimes z+e_2,\\
 p_2x&=y\otimes\psi\otimes r+f_2,&
 p_3x&=-y\otimes\ell\otimes\psi+f_3,
\end{align*}
where $\norm{e_1},\norm{e_2}\le\epsilon$ and $\norm{f_2},\norm{f_3}\le\zeta$.
Comparing the two expressions for $p_2x$ gives
$\norm{\ell\otimes z+y\otimes r}\le\epsilon+\zeta$.
Let $\alpha=\langle r,z\rangle$. Contracting with $r$ on the right gives $\norm{y+\alpha\ell}\le\epsilon+\zeta$; projecting on $r^\perp$ gives $\norm{z-\alpha r}\le\epsilon+\zeta$. Consequently
\begin{align*}
 p_1x&=\alpha\psi\otimes r\otimes r+E_1,\\
 p_2x&=-\alpha\ell\otimes\psi\otimes r+E_2,\\
 p_3x&=\alpha\ell\otimes\ell\otimes\psi+E_3,
\end{align*}
with $\norm{E_1},\norm{E_2}\le2\epsilon+\zeta\le\sqrt5R_0$ and
$\norm{E_3}\le\epsilon+2\zeta\le\sqrt5R_0$.
Since $2\norm{p_1p_3x}^2\le E$, applying $p_1$ to the third expression and using \eqref{eq:scalars} yields
\[
 |\alpha|\sqrt{\tau/2}\le\sqrt{E/2}+\sqrt5R_0
 \le(1/2+\sqrt5)R_0,
\]
where $\delta\le1/2$ was used. Thus $|\alpha|\le4R_0/\sqrt\tau$.
Minkowski's inequality in the direct sum of the three projected spaces gives
\[
 \sqrt{\langle x,hx\rangle}
 \le\sqrt3|\alpha|+\sqrt{15}R_0
 \le(4\sqrt3+\sqrt{15})R_0/\sqrt\tau.
\]
Because $(4\sqrt3+\sqrt{15})^2=63+24\sqrt5<128$, this proves \eqref{eq:Gquant}. The estimates also cover $E=0$ directly.
\end{proof}

\subsection{Completion of the classification}
If $\tau>0$, apply Lemmas~\ref{lem:quantitative} and \ref{lem:saturation}, noting that $\tau\delta/128\le1/256<\eta$. This gives exactly \eqref{eq:generalbound}. If $\tau=0$, then $c=0$ in \eqref{eq:Uframe}; unitarity forces $b=0$ and $|a|=|d|=1$. Absorb these phases into $u,v$ to obtain \eqref{eq:balancedmarker}. Theorem~\ref{thm:marker} then gives \eqref{eq:diffusive}.

If $L=R$, restrict to this common two-dimensional space. The homogeneous complex quadratic $\langle w,w|\psi\rangle$ has a nonzero zero $w$: this is the elementary fact that every binary complex quadratic is isotropic, with the identically zero polynomial included. Normalise $w$ and complete it by an orthogonal unit vector $u$. In the frame $(u,w)$ on both factors, the scaled coefficient matrix is unitary. Its $ww$ entry vanishes, so its $uu$ entry also vanishes and the two off-diagonal entries have modulus one. The forbidden state is a balanced marker with $v$ a phase multiple of $u$. The exact marker theorem applies on the whole on-site space, including unused levels.

When the intersection is one-dimensional, $L+R$ has dimension three, irrespective of the ambient $d$. The quantitative proof was given using only vectors and operators supported on this span. Equivalently, orthogonal spectator levels split the chain into open active segments and cannot lower a uniform bound. This proves every case of Theorem~\ref{thm:classification}. It proves the if-and-only-if statement without assuming that a low-energy vector in a selected sector is orthogonal to an unspecified ground space.

\section{The complete two- and three-site isospectral fibre}\label{sec:isospectral}
\begin{theorem}[Normal form and gap classification of the fibre]\label{thm:fibre}
Suppose $d=3$, the Schmidt probabilities are $(1/2,1/2,0)$, and
\begin{equation}\label{eq:h3spec}
 \spec H_3=\{0^{(21)},(1/2)^{(1)},1^{(4)},(3/2)^{(1)}\}.
\end{equation}
Up to an identical on-site unitary and an irrelevant phase, there is a unique $\theta\in[0,\pi/2]$ for which
\begin{equation}\label{eq:isopath}
 \psi=\psi_\theta=
 \frac{\cos\theta\ket{00}+\sin\theta\ket{01}-\ket{12}}{\sqrt2}.
\end{equation}
Every point has \eqref{eq:h3spec} and $\spec H_2=\{0^{(8)},1\}$. All $\theta<\pi/2$ are uniformly gapped, with
\begin{equation}\label{eq:fibrebound}
 \gamma_N(\psi_\theta)\ge\frac{\cos^2\theta}{6}\quad(N\ge2),
\end{equation}
whereas $\gamma_N(\psi_{\pi/2})=1-\cos(\pi/N)$. The bound in \eqref{eq:fibrebound} is a lower bound, not an exact formula away from the endpoint.
\end{theorem}
\begin{proof}[Normal form and short spectra]
The positive eigenvalues of $H_3$ are $1\pm s_j(V^*W)$. Since $s_j\le1/2$, none disappears at zero. Therefore \eqref{eq:h3spec} forces overlap singular values $(1/2,0,0)$ and hence $t=0$. The supports intersect in a line but are not equal, and $u,w,v$ in \eqref{eq:Uframe} form an orthonormal basis. Their independent phases, together with the phase of $\psi$, reduce the unitary coefficient matrix to
\[
 \begin{pmatrix}-s&-c\\c&-s\end{pmatrix},\qquad
 c=\sqrt\tau,\quad s=\sqrt{1-\tau}.
\]
For nonzero entries, the only phase constraint is the unitary identity
$\arg a+\arg d-\arg b-\arg c=\pi\pmod{2\pi}$, which is exactly the constraint of this real representative. Vanishing endpoint entries cause no obstruction. Set $e_0=cw-su$, $e_1=sw+cu$, $e_2=v$. This gives \eqref{eq:isopath} with $c=\cos\theta$, $s=\sin\theta$. The invariant $\tau$ determines $\theta$ uniquely.

Conversely
\[
 M_\theta=\frac1{\sqrt2}\begin{pmatrix}c&s&0\\0&0&-1\\0&0&0\end{pmatrix},\quad
 M_\theta^2=\frac12\begin{pmatrix}c^2&cs&-s\\0&0&0\\0&0&0\end{pmatrix}.
\]
The rows of $M_\theta$ have squared norm $1/2$ and are orthogonal. The sole nonzero singular value of $M_\theta^2$ is $1/2$, using $c^4+c^2s^2+s^2=1$. This proves the stated Schmidt and short-chain spectra. Gappedness except at $c=0$ already follows from Theorem~\ref{thm:classification}. The sharper bound is proved next.
\end{proof}

\subsection{A three-dimensional certificate inside the four-site range}
\begin{lemma}[Full range-Gram decomposition]\label{lem:gram}
For the fibre normal form, the operator $K-D$ defined below is unitarily
equivalent to $S(c)$, the scalars $-1/2,1/2$, four copies of
$\left(\begin{smallmatrix}1/2&1/2\\1/2&1\end{smallmatrix}\right)$,
ten scalar blocks $1/2$, and four scalar blocks $1$.
Its negative eigenspace is the single direction $(\psi,0,-\psi)$,
which lies in $\ker C$. In particular the entire Gram operator is not
positive: the relevant inequality is on $\ran C^*$.
\end{lemma}
\begin{proof}
Write $\ell=e_0$, $r=(c^2,cs,-s)^T$ in the computational frame of \eqref{eq:isopath}. Let
\[
 C=(V_1\;V_2\;V_3),\quad
 V_1x=\psi\otimes x,\quad V_2=I\otimes\psi\otimes I,\quad V_3x=x\otimes\psi,
\]
where every $V_i$ has nine columns. Then $H_4=CC^*$ and, with $K=C^*C$,
\[
 \mathcal G=C(K-D)C^*,\qquad D=\diag(\tfrac12 I_9,0_9,\tfrac12 I_9).
\]
The off-diagonal Gram blocks are
\[
 K_{12}=\tfrac12(\ket r\bra\ell\otimes I),\quad
 K_{23}=\tfrac12(I\otimes\ket r\bra\ell),\quad
 K_{13}=\ket\psi\bra\psi.
\]
Set $\alpha=c/\sqrt2$ and $\beta=\sqrt{1-c^2/2}$. A reducing five-dimensional core is spanned, in the first, middle and last column spaces respectively, by
\[
 a=r\otimes r,\ a'=(\psi-\alpha a)/\beta;\qquad
 b=\ell\otimes r;\qquad
 d=\ell\otimes\ell,\ d'=(\psi-\alpha d)/\beta.
\]
These are orthonormal within their column spaces. Reflection exchanging $(a,a')$ and $(d,d')$ splits the core. The odd block has eigenvalues $-1/2,1/2$. The negative vector is proportional to $(\psi,0,-\psi)$ and is killed by $C$. In the even basis $(a+d)/\sqrt2,(a'+d')/\sqrt2,b$, the block is
\begin{equation}\label{eq:Score}
 S(c)=\begin{pmatrix}
 (1+c^2)/2&c\sqrt{2-c^2}/2&1/\sqrt2\\
 c\sqrt{2-c^2}/2&(3-c^2)/2&0\\
 1/\sqrt2&0&1
 \end{pmatrix}.
\end{equation}
Outside this core, there are four two-dimensional blocks with diagonal $(1/2,1)$ and off-diagonal $1/2$, ten scalar blocks $1/2$, and four scalar blocks $1$. To see the four paired blocks explicitly, use the two dimensions of $\ell\otimes r^\perp$ in the middle space and their first-space partners $r\otimes r^\perp$, and the two dimensions of $\ell^\perp\otimes r$ and their last-space partners $\ell^\perp\otimes\ell$. The remaining spaces have no overlap coupling. Appendix~\ref{app:basis} specifies the complementary bases and the transformation without excluding either endpoint.
\end{proof}

\begin{proof}[Completion of the fibre bound]

For $c>0$, the leading principal minors of $S(c)$ are $(1+c^2)/2$, $3/4$, and $c^2/4$, so $S(c)>0$. Its characteristic polynomial is
\begin{equation}\label{eq:Scubic}
 \det(xI-S(c))=x(x-3/2)^2-c^2/4.
\end{equation}
If $\lambda_1\le\lambda_2\le\lambda_3$ are its eigenvalues, the sum of their pairwise products is $9/4$. Hence
$\lambda_1=(c^2/4)/(\lambda_2\lambda_3)\ge c^2/9$.
The other nonnegative blocks have least eigenvalue at least $(3-\sqrt5)/4>1/9$. Since $\ran C^*$ is orthogonal to the negative vector, we have
\[
 \mathcal G\ge(c^2/9)H_4.
\]
Lemma~\ref{lem:saturation} proves \eqref{eq:fibrebound}. At $c=0$ the marker theorem gives the exact endpoint gap. This completes Theorem~\ref{thm:fibre}.
\end{proof}

\subsection{Constant degeneracy and a dimension-minimal separation}
The recurrence itself is the $d=3,r=1$ case of equation~(11) in
Section~II of Movassagh et al.~\cite{Mov}. What is proved here is its
exact validity at every point of this specified fibre, not just at generic
interactions, together with its failure to determine gappedness.
\begin{proposition}[The ground-space dimension stays fixed along the fibre]\label{prop:degeneracy}
For every $\theta\in[0,\pi/2]$, let $D_N=\dim\ker H_N(\psi_\theta)$ and $D_0=1$. Then
\[
 D_1=3,\qquad D_N=3D_{N-1}-D_{N-2}=F_{2N+2},
\]
where $F_0=0,F_1=1$ are the Fibonacci numbers.
\end{proposition}
\begin{proof}
The orthogonal complement of the ground space is the degree-$N$ part of the two-sided tensor ideal generated by $c\,00+s\,01-12$. Orient its sole linear rewriting rule as $12\mapsto c\,00+s\,01$. Degree-lexicographic order with $1>0$ makes every replacement decrease the word. The leading word $12$ has no proper self-overlap. Disjoint replacements commute, so induction in this finite order makes the normal form well-defined and unique. Thus words avoiding $12$ form a basis of the quotient: a linear combination of such words cannot reduce to zero unless all coefficients vanish. There are $3D_{N-1}$ extensions of admissible words, and exactly $D_{N-2}$ forbidden extensions ending in $12$, because every admissible prefix can be followed by $1$. This gives the recurrence and its stated solution.
\end{proof}
The gapped and gapless endpoints therefore have the same ground-space dimension at every length, as well as the same specified short-chain spectra and Schmidt probabilities. No assertion is made that their full spectra agree beyond length three.

Put $\psi_{\mathrm d}=\psi_0$, $\psi_{\mathrm m}=\psi_{\pi/2}$ and $P_{\mathrm d,\mathrm m}=\ket{\psi_{\mathrm d,\mathrm m}}\bra{\psi_{\mathrm d,\mathrm m}}$. The sharper dimer estimate in Section~\ref{sec:dimer} gives
\begin{equation}\label{eq:mainbounds}
 \gamma_N(P_{\mathrm d})\ge1/4,\qquad
 \gamma_N(P_{\mathrm m})=1-\cos(\pi/N).
\end{equation}
The union of the left and right supports spans all three local levels at either endpoint. Together with the qubit part of Theorem~\ref{thm:classification}, this proves Corollary~\ref{cor:minimal}.

\subsection{The weaker qubit obstruction from the supplied review}\label{sec:qubitcomparison}
The unrestricted statement that Schmidt data and two short spectra cannot classify gaps is already a short consequence of Bravyi--Gosset. The following explicit pair was derived in the supplied review, rather than quoted as a pair printed in \cite{BG}:
\[
 \chi_c=\frac{\ket{00}}{\sqrt2}+\frac{\ket{01}-\ket{10}}2,
 \qquad
 \chi_g=\frac{\ket{00}}2+\frac{1+\sqrt5}{4}\ket{01}
       +i\frac{\sqrt5-1}{4}\ket{10}.
\]
Their Schmidt probabilities are $(2\pm\sqrt3)/4$, and their common three-site characteristic polynomial is
\[
 x^4\bigl((x-1)^4-\tfrac38(x-1)^2+1/256\bigr).
\]
For $T_\chi=\left(\begin{smallmatrix}\overline{\chi_{01}}&\overline{\chi_{11}}\\-\overline{\chi_{00}}&-\overline{\chi_{10}}\end{smallmatrix}\right)$,
the first state has eigenvalue moduli $1/2,1/2$ and the second has $(\sqrt5+1)/4,(\sqrt5-1)/4$. Theorem 1 of \cite{BG} makes the first gapless and the second gapped. Adding a spectator level yields a qutrit comparison. This does not have flat Schmidt probabilities. We retain the supplied exact verifier and its derivation in the evidence bundle and do not claim this weaker obstruction as an independent new theorem.
\section{The dimer family is uniformly gapped}\label{sec:dimer}
\begin{theorem}[All positive dimer weights]\label{thm:dimer}
For $a,b>0$, $a^2+b^2=1$, and
\begin{equation}\label{eq:dimer-family}
 \psi=a\ket{00}-b\ket{12},
\end{equation}
one has $\gamma_N(\ket\psi\bra\psi)\ge b^4$ for every $N\ge2$.
At either endpoint $a=0$ or $b=0$ the commuting-projector chain has gap one.
\end{theorem}

\subsection{Reducing sectors as hard-core configurations}
The only off-diagonal move is $00\leftrightarrow12$. Any symbol $1$ or $2$ not belonging to an adjacent $12$ pair is fixed under all moves: no move can act on it or produce a new $12$ pair across it. These frozen symbols split the chain into intervals. Each remaining interval is tiled by monomers $0$ and oriented dimers $12$. All such tilings of the interval lie in one connected component: delete all dimers and then insert the desired ones. Conversely no move leaves this component.

On an interval of $\ell$ physical sites, a tiling is an independent set $\eta$ on the path of $m=\ell-1$ possible dimer edges. Dimers cannot occupy adjacent edges. The positive ground amplitude is
\[
 g(\eta)=\left(\frac ab\right)^{|\eta|}.
\]
Conjugating the sector Hamiltonian by $G=\diag g$ gives the positive generator $\mathcal L=G^{-1}H G$ with insertion rate $p=a^2$ and deletion rate $d=b^2$. These are exactly rate-one single-site heat-bath updates for the hard-core law
\[
 \pi(\eta)=Z^{-1}\lambda^{|\eta|},\qquad \lambda=a^2/b^2,
 \qquad p=\lambda/(1+\lambda),\quad d=1/(1+\lambda).
\]
Insertion is allowed only when the neighbouring hard-core sites are empty. In particular $\mathcal L$ is reversible in $L^2(\pi)$ and has the same spectrum as the quantum sector Hamiltonian. The full chain is a direct sum of tensor sums of these interval Hamiltonians and zero one-dimensional sectors. It remains to bound every finite hard-core path uniformly, including arbitrary fugacity.

\subsection{A two-site heat-bath comparison}
\begin{lemma}[Hard-core path bound]\label{lem:hardcore}
For every path length $m\ge1$ and every fugacity $\lambda>0$, the rate-one single-site heat-bath generator satisfies
\begin{equation}\label{eq:hardcore-gap}
 \gap(\mathcal L)\ge(1+\lambda)^{-2}=d^2.
\end{equation}
\end{lemma}
\begin{proof}
Let $E_i$ denote conditional expectation resampling site $i$, so
$\mathcal L=\sum_i(I-E_i)$. Consider the block heat-bath generator
\[
 \mathcal B=(I-E_{\{1\}})+\sum_{i=1}^{m-1}(I-E_{\{i,i+1\}})
                 +(I-E_{\{m\}}).
\]
For $m=1$, retain both copies of the singleton. Every site belongs to exactly two blocks.

First we show $\gap(\mathcal B)\ge2d$. Equip independent sets with Hamming distance. This distance is geodesic for single-site admissible changes: remove the occupations absent from the target and then add its missing occupations. Consider two configurations differing only at site $j$. Each of the two blocks containing $j$ has identical outside conditions for the two configurations and can be resampled identically, eliminating the discrepancy. At most two other blocks can feel the changed boundary site. Each creates expected additional Hamming distance at most $p$ under a suitable coupling.

Here is the complete nontrivial two-site coupling calculation. With neither site blocked, the conditional weights of $00,10,01$ are $1,\lambda,\lambda$. If the exterior discrepancy blocks the first site, the other conditional weights are $1,\lambda$ on $00,01$. Couple the common mass at $00$ and $01$ identically. The remaining mass comes from $10$: send
\[
 \frac{\lambda}{(1+\lambda)(1+2\lambda)}\ \text{to }00,
 \qquad
 \frac{\lambda^2}{(1+\lambda)(1+2\lambda)}\ \text{to }01.
\]
The transport distances are one and two, so the expected new distance is
$\lambda/(1+\lambda)=p$. If the other site is blocked by the unchanged exterior condition, the calculation is a single Bernoulli update, again costing at most $p$. Singleton blocks have the same bound.

Synchronising the rate-one clocks therefore gives infinitesimal expected-distance drift at most $-2+2p=-2d$ for neighbouring configurations. Concatenating the couplings along geodesics extends this bound to arbitrary configurations, and iteration gives contraction by $e^{-2dt}$ for the block semigroup in the Hamming Wasserstein distance. Equivalently it contracts the Lipschitz seminorm of functions by this factor. Every nonconstant eigenfunction has nonzero finite Lipschitz seminorm; applying the contraction to an eigenfunction proves that every nonzero eigenvalue of $\mathcal B$ is at least $2d$.

Next compare block and single-site Dirichlet forms. Conditional on the exterior of a free pair, the positive generator $\mathcal L_i+\mathcal L_{i+1}$ on $00,10,01$ is
\[
 \begin{pmatrix}2p&-p&-p\\-d&d&0\\-d&0&d\end{pmatrix},
\]
with eigenvalues $0,d,1+p$. If an exterior neighbour blocks a site, the conditional nonzero gap is one, or the space is one-dimensional. Hence in every exterior condition
\[
 I-E_{\{i,i+1\}}\ \le\ d^{-1}(\mathcal L_i+\mathcal L_{i+1})
\]
as an $L^2(\pi)$ quadratic-form inequality. The same bound applies to singleton blocks. Summing, using the two-fold incidence of every site, yields
$\mathcal B\le2d^{-1}\mathcal L$. For every mean-zero function $f$,
\[
 \langle f,\mathcal Lf\rangle_\pi
 \ge\frac d2\langle f,\mathcal Bf\rangle_\pi
 \ge d^2\norm f_\pi^2.
\]
This is \eqref{eq:hardcore-gap}.
\end{proof}
\begin{proof}[Proof of Theorem~\ref{thm:dimer}]
Apply Lemma~\ref{lem:hardcore} to every nontrivial interval factor. The positive spectrum of a tensor sum is bounded below by the minimum factor gap, and direct sums preserve the common lower bound. Thus every positive eigenvalue of the full quantum chain is at least $d^2=b^4$. At $a=0$ or $b=0$ all local projectors are diagonal and the smallest positive energy is one.
\end{proof}
In particular $a=b=1/\sqrt2$ gives the first inequality in \eqref{eq:mainbounds}. Section~\ref{sec:certificate} supplies an independent, weaker $1/10$ bound at this point using only a rational four-site calculation. The new saturation proof gives a third, independent lower bound. The heat-bath estimate is retained because it covers every positive dimer weight, including unequal Schmidt probabilities.

\section{A rational four-site certificate and perturbation stability}\label{sec:certificate}
\begin{theorem}[Certified open region]\label{thm:ball}
For every rank-one orthogonal projector $Q$ with
$\epsilon=\norm{Q-\Pd}<1/45$, one has
\begin{equation}\label{eq:ball}
 \inf_{N\ge2}\gamma_N(Q)\ge\frac1{10}-\frac92\epsilon>0.
\end{equation}
In particular the closed radius-$1/90$ ball has gap at least $1/20$.
The full-Schmidt-rank vector
\begin{equation}\label{eq:star}
 \psi_* =\frac{101\ket{00}-100\ket{12}+\ket{21}}{\sqrt{20202}}
\end{equation}
satisfies this bound, although its three-site gap is strictly below $1/2$.
\end{theorem}
This is a neighbourhood in the full complex rank-one projector space, not a neighbourhood confined to a chosen real slice. The local projector rank remains one; the vector's Schmidt rank is a different notion.

\subsection{An open-chain finite-size inequality}
\begin{lemma}[Clipped-window criterion]\label{lem:windows}
Fix $m\ge3$ and write $\delta_m=\min_{2\le k\le m}\gamma_k(P)$. For $N\ge m$,
\begin{equation}\label{eq:finite-size}
 \gamma_N(P)\ge\frac{(m-1)\delta_m-1}{m-2}
\end{equation}
whenever the right-hand side is positive.
\end{lemma}
\begin{proof}
Write $p_i=P_{i,i+1}$ for $1\le i\le N-1$, extending $p_i=0$ outside this range. For each integer starting point $t$ whose window meets the chain, set
$B_t=\sum_{i=t}^{t+m-2}p_i$. This includes clipped windows at both ends, so
$\sum_tB_t=(m-1)H_N$. Every nonzero window is a shorter open-chain Hamiltonian on between two and $m$ sites; therefore
\[
 \mathcal A:=\sum_tB_t^2\ge(m-1)\delta_m H_N.
\]
Expanding exactly gives
\[
 \mathcal A=(m-1)H_N+
 \sum_{\substack{i<j\\j-i\le m-2}}(m-1-j+i)\{p_i,p_j\}.
\]
The coefficient of each adjacent anticommutator is $m-2$. Every other anticommutator is nonnegative, since its projectors have disjoint support. All its coefficients are at most $m-2$. Hence
$\mathcal A\le(m-2)H_N^2+H_N$.
Combining these inequalities and using the spectral theorem proves the result. The shorter chains $N<m$ are handled by their individual gaps, not by assuming periodic boundaries.
\end{proof}

\subsection{The exact five-state calculation}
For $P=\Pd$, the computational-basis components under $00\leftrightarrow12$ on four sites have sizes and multiplicities
\[
 1^{(36)},\quad2^{(14)},\quad3^{(4)},\quad5^{(1)}.
\]
The singleton blocks vanish; the two-state blocks have spectrum $0,1$; the three-state blocks have spectrum $0,1/2,3/2$. The only remaining block, in the order
$0000,0120,1200,1212,0012$, is
\begin{equation}\label{eq:K}
 K=\frac12\begin{pmatrix}
 3&-1&-1&0&-1\\-1&1&0&0&0\\-1&0&2&-1&0\\
 0&0&-1&2&-1\\-1&0&0&-1&2
 \end{pmatrix}.
\end{equation}
Let $B=K^2-\tfrac25K$ and let $C$ be the leading four-by-four principal submatrix of $B$. The exact identity
\begin{equation}\label{eq:LDL}
 C=LDL^T,\qquad
 L=\begin{pmatrix}
 1&0&0&0\\-1/3&1&0&0\\-7/16&-3&1&0\\
 5/24&5&-26/109&1
 \end{pmatrix},\qquad
 D=\diag\left(\frac{12}5,\frac1{30},\frac{109}{320},\frac{78}{545}\right)
\end{equation}
has strictly positive diagonal pivots. Since $B$ has zero row sums,
$B=TCT^T$ for $T=[I_4;-\mathbf1^T]$, and hence $B\ge0$. This proves that the nonzero spectrum of $K$ is at least $2/5$. The finite block classification is directly checkable by the listed move rule; the verifier also constructs the full $81\times81$ matrix independently by tensor products and checks all component boundaries.

We have consequently established the exact bounds
\begin{equation}\label{eq:localbounds}
 \gamma_2=1,\qquad\gamma_3=1/2,\qquad\gamma_4\ge2/5.
\end{equation}
As a cross-check, the nontrivial block's characteristic polynomial is
\[
 \det(xI-K)=\frac14x(x-1)(4x^3-16x^2+18x-5).
\]
Its least positive root is approximately $0.4149567567$; this decimal is not used in the proof. Lemma~\ref{lem:windows}, with $m=4$, gives a uniform gap at least $1/10$.

\subsection{Why degeneracy cannot invalidate the perturbation bound}
A gap bound at a highly degenerate frustration-free point cannot in general be perturbed merely by applying Weyl's inequality: new small positive eigenvalues might emerge from its kernel. Here local rank maxima prevent that failure.

\begin{lemma}[Short-chain rank maximality]\label{lem:rankmax}
For every rank-one qutrit projector $Q$,
\begin{equation}\label{eq:ranks}
 \rank H_2(Q)\le1,\qquad \rank H_3(Q)\le6,\qquad \rank H_4(Q)\le26.
\end{equation}
The dimer projector attains all three maxima.
\end{lemma}
\begin{proof}
The first two inequalities follow from the individual range dimensions. On four sites each bond-projector range has dimension nine. The first and third share the one-dimensional product range of $Q_{12}Q_{34}$, so their total span has dimension at most $9+9+9-1=26$. The dimer block decomposition gives kernel dimensions $8,21,55$, attaining the maxima.
\end{proof}

For $\epsilon=\norm{Q-\Pd}$,
\[
 \norm{H_k(Q)-H_k(\Pd)}\le(k-1)\epsilon.
\]
Weyl's inequality keeps all the existing positive eigenvalues positive at the stated small radii. The universal upper bounds \eqref{eq:ranks} prohibit any additional positive eigenvalues. Thus the kernel dimensions stay fixed for $k\le4$, and
\begin{equation}\label{eq:pertlocal}
 \gamma_2(Q)=1,\qquad
 \gamma_3(Q)\ge\frac12-2\epsilon,\qquad
 \gamma_4(Q)\ge\frac25-3\epsilon.
\end{equation}
For $\epsilon<1/45$, the last lower bound is the smallest. Substitution into \eqref{eq:finite-size} proves \eqref{eq:ball}; the shorter chains satisfy the same global bound directly from \eqref{eq:pertlocal}.

For \eqref{eq:star}, the unnormalised coefficient matrix is
\[
 C_* =\begin{pmatrix}101&0&0\\0&0&-100\\0&1&0\end{pmatrix},
 \qquad \det C_*=10100\ne0.
\]
The exact squared distance of rank-one projectors is
\[
 \norm{\ket{\psi_*}\bra{\psi_*}-\Pd}^2
 =1-|\langle\psi_{\mathrm d},\psi_*\rangle|^2
 =\frac1{13468}<\frac1{90^2}.
\]
On the other hand,
\[
 \norm{\overline M_*M_*}=\frac{10201}{20202}>\frac12,
 \qquad \gamma_3(\ket{\psi_*}\bra{\psi_*})=\frac{10001}{20202}<\frac12.
\]
The three-site overlap test therefore gives no positive thermodynamic lower bound, while our four-site certificate proves $\inf_N\gamma_N\ge1/20$. This completes Theorem~\ref{thm:ball}.

\section{Rank two: conserved order and exponential bottlenecks}\label{sec:ranktwo}
\begin{theorem}[Exponential gap upper bound for opposing species]\label{thm:ranktwo-main}
For $0<a<1<b$, let
\begin{equation}\label{eq:ranktwo}
 \phi_1=\frac{\ket{10}-a\ket{01}}{\sqrt{1+a^2}},\qquad
 \phi_2=\frac{\ket{20}-b\ket{02}}{\sqrt{1+b^2}},\qquad
 P(a,b)=\sum_{s=1}^2\ket{\phi_s}\bra{\phi_s}.
\end{equation}
The two vectors are orthonormal, so $P(a,b)$ is a rank-two projector. Put
$A=-\log a$ and $B=\log b$. For every $N\ge 2+\lceil B/A\rceil$,
\begin{equation}\label{eq:ranktwo-exp}
 \gamma_N(P(a,b))\le
 \frac{1+(b/a)^2}{1+a^2}
 \exp\left[-\frac{2AB}{A+B}(N-1)\right].
\end{equation}
Each chain with only one of the two rank-one constituents is uniformly gapped, with gap infimum $1-2q/(1+q^2)>0$, where $q=a$ or $q=b$.
For the rational projectors defined by
\begin{equation}\label{eq:rational-two}
 \phi_1=\frac{2\ket{10}-\ket{01}}{\sqrt5},\qquad
 \phi_2=\frac{\ket{20}-2\ket{02}}{\sqrt5},
\end{equation}
one has the sharper simple bound
\begin{equation}\label{eq:rational-exp}
 \gamma_{2L+1}\le\frac85\,4^{-L}\qquad(L\ge1),
\end{equation}
whereas both constituent families have gap infimum $1/5$.
\end{theorem}
\subsection{Exact one-vacancy reduction}
In \eqref{eq:ranktwo}, $0$ is a vacancy and $1,2$ are particles which never exchange with one another. The order of the nonzero symbols is conserved. Fix any word $\sigma=\sigma_1\cdots\sigma_m$ over $\{1,2\}$ and take $N=m+1$. Its one-vacancy sector has the orthonormal basis
\[
 e_j=\ket{\sigma_1\cdots\sigma_j\,0\,\sigma_{j+1}\cdots\sigma_m},
 \qquad0\le j\le m.
\]
It is a reducing subspace, and the restricted Hamiltonian is $D^*D$, with row $i$ of $D$ equal to
\begin{equation}\label{eq:vacancyD}
 \frac{-q_{\sigma_i}e_{i-1}^T+e_i^T}{\sqrt{1+q_{\sigma_i}^2}},
 \qquad q_1=a,\quad q_2=b.
\end{equation}
The kernel in this sector is exactly one-dimensional, spanned by the strictly positive vector
\begin{equation}\label{eq:g}
 g_0=1,\qquad g_j=\prod_{i=1}^j q_{\sigma_i}.
\end{equation}
This identifies the ground-state component explicitly. A trial vector in this sector orthogonal to $g$ is therefore orthogonal to the \emph{entire} ground space of the full chain, not just to a selected product ground state.

\subsection{Opposing biases produce a two-well trial state}
Assume $0<a<1<b$ and let $A=-\log a$, $B=\log b$. For $m=N-1\ge1+\lceil B/A\rceil$ choose
\[
 k=\left\lceil\frac{mB}{A+B}\right\rceil,\qquad
 \sigma=1^k2^{m-k},
\]
with $1\le k\le m-1$. Then $g_j=a^j$ up to $j=k$, and
$g_j=a^k b^{j-k}$ afterwards. In particular $g_0=1$, $a/b\le g_m\le1$, and $g_k=a^k$ is exponentially small. The inhomogeneous effective potential comes from a conserved particle order; the local quantum Hamiltonian itself remains translation invariant.

Put
\[
 W_L=\sum_{j=0}^{k-1}g_j^2,\qquad W_R=\sum_{j=k}^{m}g_j^2,
 \qquad
 f_j=\begin{cases}W_Rg_j,&j<k,\\-W_Lg_j,&j\ge k.\end{cases}
\]
Then $\langle f,g\rangle=0$. Every row of $Df$ vanishes except the row crossing the cut. Exact cancellation gives
\begin{equation}\label{eq:rayleigh}
 \frac{\langle f,H_N f\rangle}{\langle f,f\rangle}
 =\frac{a^{2k}}{1+a^2}\left(\frac1{W_L}+\frac1{W_R}\right).
\end{equation}
Since $W_L\ge1$, $W_R\ge(a/b)^2$, and
$a^{2k}\le\exp[-2ABm/(A+B)]$, the variational principle proves \eqref{eq:ranktwo-exp}. Exchanging the species proves the same conclusion in the other opposing-bias quadrant.

For $a=1/2$, $b=2$, $m=2L$, take $k=L$. Both endpoints of $g$ equal one, so $W_L,W_R\ge1$. Equation~\eqref{eq:rayleigh} gives \eqref{eq:rational-exp} directly. The verifier checks the exact rational trial norm, orthogonality, energy, and bound for $1\le L\le50$; the argument above proves the formula for every $L$.

\subsection{Individually gapped constituents}
For the Hamiltonian with only $\ket{\phi_1}\bra{\phi_1}$, every occurrence of symbol $2$ is a fixed spectator. Fixing its positions splits the chain into open qubit segments on $\{0,1\}$. Each segment is the $u=v$ case of Theorem~\ref{thm:marker}, with gap
\[
 1-\frac{2a}{1+a^2}\cos(\pi/\ell)
\]
for segment length $\ell\ge2$. The active sector with no spectators attains the full-chain gap, and its infimum is $1-2a/(1+a^2)$. The same reasoning applies to the second constituent. This proves the remaining statements of Theorem~\ref{thm:ranktwo-main}. It also shows why adding positive, individually gapped frustration-free summands is not enough: their common ground space can have very small principal-angle excitations.

\subsection{A complementary certified gapped region}
For completeness the same rank-two family admits an elementary gapped region. Write
$q_- =\min(a,b)$, $q_+=\max(a,b)$. A direct six-dimensional overlap calculation gives
\begin{equation}\label{eq:two-overlap}
 c(a,b)=\norm{P(a,b)_{12}P(a,b)_{23}}
       =\frac{q_+}{\sqrt{(1+q_-^2)(1+q_+^2)}}.
\end{equation}
Indeed the squared singular values of the overlap of its two range isometries are
\[
 \frac{a^2}{(1+a^2)^2}\ (\text{twice}),\quad
 \frac{b^2}{(1+b^2)^2}\ (\text{twice}),\quad
 \frac{a^2}{(1+a^2)(1+b^2)},\quad
 \frac{b^2}{(1+a^2)(1+b^2)}.
\]
The largest is the square of \eqref{eq:two-overlap}. Therefore
\begin{equation}\label{eq:rank2region}
 4q_+^2<(1+q_-^2)(1+q_+^2)
 \quad\Longrightarrow\quad
 \inf_N\gamma_N(P(a,b))\ge1-2c(a,b)>0.
\end{equation}
The opposing-bias region is open within this two-parameter family; stability under arbitrary rank-two perturbations is not proved. For equal biases $a=b=q$, the lower bound of Lemma~\ref{lem:overlap} and the one-species upper bound coincide, yielding the exact formula
\[
 \gamma_N(P(q,q))=1-\frac{2q}{1+q^2}\cos(\pi/N).
\]
If either bias equals one, a one-species sector instead gives
$\gamma_N\le1-\cos(\pi/N)$, so those lines are gapless. We do \emph{not} classify the remaining same-direction parameter regions outside \eqref{eq:rank2region}.

\begin{remark}[Not the full PVBS model]\label{rem:pvbs}
The hopping-only projector \eqref{eq:ranktwo} lacks both same-species exclusion penalties and unlike-species interchange terms present in the two-species Product Vacua and Boundary States models. The PVBS conclusion that adding a species does not introduce a new gapless phase \cite{Bishop} therefore does not apply to it. Our exponentially small gaps do not contradict that result.
\end{remark}


\subsection{Definition-level Motzkin and free-Motzkin comparisons}\label{sec:motzkin}
Write $P_{\mathrm d}=\ket{00-12}\bra{00-12}/2$. In the notation $u=1,d=2$, the original Motzkin bulk projector is exactly
\[
 P_{\mathrm{Motzkin}}=P_{\mathrm d}+P(1,1).
\]
Bravyi et al. add endpoint penalties $\ket2\bra2_1+\ket1\bra1_N$ to select the unique Motzkin ground state \cite{Motzkin}. In Step~1 of their proof, immediately before equation~(6), they separate the hopping and pair-creation terms and explicitly use the ferromagnetic hopping gap $1-\cos(\pi/N)$. Thus neither the unbiased hopping interaction nor that gap formula is claimed as new here.

Salberger, Padmanabhan and Korepin study the hopping-only free-Motzkin model with periodic boundaries \cite{FreeMotzkin}. In their explicit equations (2.2) and (2.5), each local difference-vector operator is \emph{twice} its normalised rank-one projector. Their bond interaction is therefore $2P(1,1)$, despite their use of the word projector. Their displayed relation $\widehat e_j^2=2\widehat e_j$ confirms this normalisation. It must be divided by two before comparison with our energy units, and a periodic result is not silently substituted for an open-chain result.

For the area-weighted model, denote the area weight by $s>0$ to distinguish it from the support angle $t$. The forbidden vectors used by Andrei, Lemm and Movassagh are, up to signs,
\[
 \frac{s\ket{01}-\ket{10}}{\sqrt{1+s^2}},\qquad
 \frac{\ket{02}-s\ket{20}}{\sqrt{1+s^2}},\qquad
 \frac{\ket{12}-s\ket{00}}{\sqrt{1+s^2}}.
\]
Thus the normalised bulk projector is
\begin{equation}\label{eq:motzkindecomposition}
 P(s,s^{-1})+\frac{\ket{12-s00}\bra{12-s00}}{1+s^2}.
\end{equation}
Their gap theorem for $0<s<1$ concerns the full Motzkin interaction, with pinning boundaries in its main statement; its proof also controls the open unpinned bulk chain \cite{ALM}. Our opposing-bias theorem instead concerns the first summand alone, without either the pair-creation term or pinning. Removing these terms enlarges the kernel, so the full-model theorem does not imply a hopping-only gap.

In particular the hopping-only open chain has one ground vector in each conserved particle-word sector. For a fixed word $\sigma_1\cdots\sigma_m$, any placement of its ordered particles can be reached by hopping through vacancies. The amplitudes proportional to $\prod_{j=1}^m q_{\sigma_j}^{-x_j}$, where $x_j$ is the position of particle $j$, satisfy every local zero-energy equation. Connectivity makes the sector kernel one-dimensional. Hence its total ground-space dimension is $\sum_{m=0}^N2^m=2^{N+1}-1$. Pair creation does not preserve those sectors. This identifies why the omitted term changes the mathematical problem.

\begin{center}
\small
\begin{tabular}{p{.22\textwidth}p{.29\textwidth}p{.39\textwidth}}
\toprule
Model & Local interaction and boundary & What the comparison establishes\\
\midrule
Original Motzkin \cite{Motzkin} & Rank-three bulk $P_{\mathrm d}+P(1,1)$; pinned ends & Direct antecedent of both components; the unbiased hopping reduction and its open gap are already used.\\[3pt]
Free Motzkin \cite{FreeMotzkin} & $2P(1,1)$ in the displayed convention; periodic & Direct hopping-only antecedent; divide energies by two and retain the boundary distinction.\\[3pt]
Area-weighted Motzkin \cite{ALM} & Rank-three projector \eqref{eq:motzkindecomposition}; pinning in main theorem & Full model gapped for $0<s<1$; does not classify its hopping-only component.\\[3pt]
Present opposing-bias result & Rank-two $P(a,b)$; open, unpinned & Explicit exponential upper bound for $0<a<1<b$, including the hopping part at $(s,s^{-1})$. No matching asymptotic lower bound.\\
\bottomrule
\end{tabular}
\end{center}
The state $(s00-12)/\sqrt{1+s^2}$ in the omitted term belongs to the dimer family of Section~\ref{sec:dimer}. Its independent uniform gap does not settle the gap of a sum with a different common kernel. The distinction from the full PVBS model in Remark~\ref{rem:pvbs} is similarly a distinction of local terms and ground spaces, not a contradiction of that model's theorem.

\section{Evidence, reproducibility and remaining questions}\label{sec:evidence}
\subsection{Proof dependencies and checks}
The principal dichotomy follows from the geometric overlap identity, the simple extremal pair mode, the quantitative four-site inequality, and the clipped-window summation. The marker sector provides an exact matching upper bound on the gapless set. The fibre normal form is a unitary-frame calculation; its stronger bound uses the three-dimensional matrix \eqref{eq:Score}. These are analytic proofs for arbitrary length and exact parameters, rather than extrapolations from diagonalisation.

The new exact verifier checks 43 groups of finite identities. It supplies seven positive rational/Hermitian $LDL^*$ certificates on the full four-site range: four general flat-Schmidt examples, including a genuinely complex one, and three fibre points. It checks the extremal-vector contractions, the displayed even core, its reducing property, the cubic polynomial and leading minors, and the clipped-window coefficient counts at lengths 2 through 30. The analytic proof, not this finite list, establishes parameter-uniform validity. The exact-input utility \texttt{classify\_flat.py} implements the support/intersection decision rule and rejects floating inputs.

A separate seeded numerical programme varies complex Schmidt-frame unitaries, nontrivial support angles and common complex on-site rotations. It also checks the marker surface, coincident supports, fibre points and four-dimensional embeddings. Its 146 full-Hamiltonian tests and 24 complex four-site relative-form tests are numerical diagnostics only; the $10^{-9}$ zero threshold is not used in a theorem. The supplied version-0.1 and referee programmes are retained with their original provenance and rerun separately. The package distinguishes historical supplied checks from newly executed checks. None of these programmes is a proof assistant.

\subsection{What the revised release establishes}
The principal result is a complete classification of a sharply specified Schmidt class, including a necessary and sufficient condition, an exact gap throughout its gapless set, and a positive analytic bound on its complement. It also gives a complete normal form and gap classification of the entire specified short-spectrum fibre. These are stronger than an isolated locally-isospectral pair and resolve the interpolation left open in version 0.1.

The arbitrary-complex rank-one neighbourhood in Section~\ref{sec:certificate} leaves the flat Schmidt class and contains full-Schmidt-rank interactions. It remains a sufficient region, not a classification of such vectors. The rank-two construction proves only an exponential \emph{upper} bound in its opposing-bias family, not the optimal exponent or the asymptotic full-chain gap. Neither the remaining same-direction rank-two parameters nor arbitrary rank-two perturbations are classified. Periodic gaps, infinite-volume GNS gaps, and the complete Bravyi--Gosset qutrit problem remain outside the claims.

The most consequential next checks are an unaffiliated adversarial reading of Lemmas~\ref{lem:saturation} and \ref{lem:quantitative}, verification of the fibre normal form against alternative conventions, and specialist prior-art review of the classification. A supplied version-1.0 review covered these arguments; its independence is not authenticated by this publication. No REF star rating follows from an internal proof audit or from the number of computed examples.

\paragraph{Executable domain.}
The classifier accepts parameter-free exact matrices only. It rejects every
free symbol before rank or null-space computation. For example,
$M(x)=\left(\begin{smallmatrix}(1-x^2)/(1+x^2)&0&2x/(1+x^2)\\0&-1&0\\0&0&0\end{smallmatrix}\right)/\sqrt2$
must first be specialised: $x=0$ is gapless, whereas $x=1$ is uniformly
gapped. Generic symbolic rank cannot certify the exceptional strata.
This restriction repairs the implementation, not the analytic theorem.

\paragraph{Declarations and availability.}
The source, exact-input classifier and finite checks accompany this manuscript.
The work is theoretical and uses no human-participant or personal data.
The scholarly creator is Anonymous. AI systems contributed mathematical
development, exposition and revision; complete originating model/run details
are not available. The publication agent performed the documented revision
and producer replay. The maintainer publishes the materials but is not thereby
credited with a research contribution. No specific research funding or
conflict declaration was supplied; neither is inferred from that absence.

\appendix
\section{An explicit endpoint-safe Gram basis}\label{app:basis}
Use the five core columns in Lemma~\ref{lem:gram}. Choose orthonormal
bases $q_1,q_2$ of $r^\perp$ and $p_1,p_2$ of $\ell^\perp$.
In the ordered direct sum of first, middle and last nine-dimensional spaces,
the four paired blocks have columns
\[
(r\otimes q_j,0,0),\ (0,\ell\otimes q_j,0);\qquad
(0,0,p_j\otimes\ell),\ (0,p_j\otimes r,0),\quad j=1,2.
\]
The four unit scalar blocks are $(0,p_i\otimes q_j,0)$.
For the ten half-unit scalar blocks, use any orthonormal bases of
\[
 (r^\perp\otimes\C^3)\cap\psi^\perp
 \quad\hbox{in the first space},\qquad
 (\C^3\otimes\ell^\perp)\cap\psi^\perp
 \quad\hbox{in the last space}.
\]
Each space has dimension five because the corresponding component of
$\psi$ has norm $\beta=\sqrt{1-c^2/2}\ge1/\sqrt2$.
These prescriptions are an explicit symbolic transformation: project the
standard tensor basis onto each displayed subspace, discard zero residuals,
and apply exact Gram--Schmidt in lexicographic order. Within the core use
the even columns displayed in the proof and the odd columns
$(a-d)/\sqrt2,(a'-d')/\sqrt2$. The odd negative and positive eigenvectors
have coordinates $(\alpha,\beta)$ and $(-\beta,\alpha)$ respectively.
All denominators used for the core are nonzero at both endpoints.
The tensor orthogonality identities and the three off-diagonal Gram blocks
give the stated block matrix by direct multiplication. This proves the
all-parameter decomposition; the accompanying finite checks remain
corroboration rather than a substitute for this argument.

\begin{thebibliography}{99}
\bibitem{BG} Bravyi, S., \& Gosset, D. (2015). Gapped and gapless phases of frustration-free spin-$1/2$ chains. \emph{Journal of Mathematical Physics, 56}, 061902. \url{https://arxiv.org/abs/1503.04035}.
\bibitem{Lemm} Lemm, M. (2019). Gaplessness is not generic for translation-invariant spin chains. \emph{Physical Review B, 100}, 035113. \url{https://doi.org/10.1103/PhysRevB.100.035113}. Version consulted: \url{https://arxiv.org/html/1903.00108v2}.
\bibitem{HJL} Hunter-Jones, N., \& Lemm, M. (2025). \emph{Two classes of quantum spin systems that are gapped on any bounded-degree graph} [Preprint]. \url{https://arxiv.org/abs/2509.22438}.
\bibitem{Rai} Rai, K. S., Kull, I., Emonts, P., Tura, J., Schuch, N., \& Baccari, F. (2026). A hierarchy of spectral gap certificates for frustration-free spin systems. \emph{Quantum, 10}, 2065. \url{https://doi.org/10.22331/q-2026-04-13-2065}.
\bibitem{Mov} Movassagh, R., Farhi, E., Goldstone, J., Nagaj, D., Osborne, T. J., \& Shor, P. W. (2010). Unfrustrated qudit chains and their ground states. \emph{Physical Review A, 82}, 012318. \url{https://arxiv.org/abs/1001.1006}.
\bibitem{Motzkin} Bravyi, S., Caha, L., Movassagh, R., Nagaj, D., \& Shor, P. W. (2012). Criticality without frustration for quantum spin-1 chains. \emph{Physical Review Letters, 109}, 207202. \url{https://doi.org/10.1103/PhysRevLett.109.207202}. \url{https://arxiv.org/abs/1203.5801}.
\bibitem{FreeMotzkin} Salberger, O., Padmanabhan, P., \& Korepin, V. (2018). \emph{Non-interacting Motzkin chain---Periodic boundary conditions} [Preprint]. \url{https://arxiv.org/abs/1809.00709}.
\bibitem{ALM} Andrei, R., Lemm, M., \& Movassagh, R. (2026). \emph{The spin-one Motzkin chain is gapped for any area weight $t<1$} [Preprint, version 2, 23 May 2026; originally submitted 2022]. \url{https://arxiv.org/abs/2204.04517}.
\bibitem{Bishop} Bishop, M. R. (2017). \emph{Spectral gaps for the two-species Product Vacua and Boundary States models on the $d$-dimensional lattice} [Preprint version consulted]. \url{https://arxiv.org/abs/1705.04755}.
\end{thebibliography}
\end{document}
